%======================================================================
% 中学试卷模板（Ayaka Launcher 内置模板库）
% 引擎：XeLaTeX；编译命令：latexmk -xelatex main.tex
% 字体：Fandol（TeX Live 自带，无需安装系统字体）
% 说明：按题型分区（选择 / 填空 / 解答），替换标题与题目即可使用。
%======================================================================
\documentclass[zihao=-4,a4paper]{ctexart}

\usepackage{amsmath,amssymb}
\usepackage[a4paper,margin=2.2cm]{geometry}
\usepackage{enumitem}
\usepackage{fancyhdr}
\usepackage{lastpage}

%---------- 页眉页脚 ----------
\pagestyle{fancy}
\fancyhf{}
\setlength{\headheight}{14pt}
\fancyhead[L]{\small 2025\,--\,2026 学年度第二学期}
\fancyhead[R]{\small 高一数学}
\fancyfoot[C]{\small 第 \thepage\ 页\quad 共 \pageref{LastPage} 页}
\renewcommand{\headrulewidth}{0.4pt}

%---------- 题目编号格式 ----------
\newlist{choice}{enumerate}{1}
\setlist[choice]{label=\textcircled{\small \Alph*}, itemsep=4pt, topsep=4pt}

\begin{document}

%---------- 试卷头 ----------
\thispagestyle{plain}
\begin{center}
  {\zihao{2}\heiti 2025\,--\,2026 学年度第二学期期末考试}\\[0.6em]
  {\zihao{3}\heiti 高一数学试卷}\\[0.8em]
  {\zihao{-4}（考试时间：120 分钟\quad 满分：150 分）}
\end{center}

\vspace{0.4em}
\noindent
\begin{tabular}{@{}p{0.31\textwidth}p{0.31\textwidth}p{0.31\textwidth}@{}}
  班级：\underline{\hspace{2.4cm}} &
  姓名：\underline{\hspace{2.4cm}} &
  考号：\underline{\hspace{2.4cm}} \\
\end{tabular}

\vspace{0.3em}
\noindent\rule{\textwidth}{0.4pt}
\noindent{\small 考生注意：本试卷共三大题。请在答题卡指定位置作答，写在试卷上的答案无效。}
\noindent\rule{\textwidth}{0.4pt}

%======================================================================
\section*{一、选择题（本大题共 8 小题，每小题 5 分，共 40 分）}

\begin{enumerate}[label=\textbf{\arabic*.}, itemsep=6pt]
  \item 已知集合 $A = \{x \mid x^2 - 3x + 2 = 0\}$，$B = \{1, 2, 3\}$，则 $A \cap B =$
  \begin{choice}
    \item $\{1\}$
    \item $\{1, 2\}$
    \item $\{2, 3\}$
    \item $\{1, 2, 3\}$
  \end{choice}

  \item 函数 $f(x) = \lg\left(x^2 - 4\right)$ 的定义域为
  \begin{choice}
    \item $(-\infty, -2)$
    \item $(2, +\infty)$
    \item $(-2, 2)$
    \item $(-\infty, -2) \cup (2, +\infty)$
  \end{choice}

  \item 已知向量 $\boldsymbol{a} = (1, 2)$，$\boldsymbol{b} = (x, -1)$，若 $\boldsymbol{a} \perp \boldsymbol{b}$，则 $x =$
  \begin{choice}
    \item $-2$
    \item $2$
    \item $-\dfrac{1}{2}$
    \item $\dfrac{1}{2}$
  \end{choice}

  \item 等差数列 $\{a_n\}$ 中，$a_2 + a_8 = 10$，则 $a_5 =$
  \begin{choice}
    \item $3$
    \item $4$
    \item $5$
    \item $6$
  \end{choice}

  \item 函数 $y = \sin\left(2x + \dfrac{\pi}{3}\right)$ 的最小正周期为
  \begin{choice}
    \item $\dfrac{\pi}{2}$
    \item $\pi$
    \item $2\pi$
    \item $4\pi$
  \end{choice}

  \item 从 5 名男生和 3 名女生中选出 3 人参加活动，则至少有 1 名女生的选法共有
  \begin{choice}
    \item $45$ 种
    \item $46$ 种
    \item $56$ 种
    \item $55$ 种
  \end{choice}

  \item 直线 $l\colon x + \sqrt{3}y - 2 = 0$ 的倾斜角为
  \begin{choice}
    \item $\dfrac{\pi}{6}$
    \item $\dfrac{\pi}{3}$
    \item $\dfrac{2\pi}{3}$
    \item $\dfrac{5\pi}{6}$
  \end{choice}

  \item 已知函数 $f(x)$ 是定义在 $\mathbf{R}$ 上的奇函数，且当 $x \geq 0$ 时 $f(x) = x^2$，则 $f(-2) =$
  \begin{choice}
    \item $-4$
    \item $4$
    \item $-2$
    \item $2$
  \end{choice}
\end{enumerate}

%======================================================================
\section*{二、填空题（本大题共 6 小题，每小题 5 分，共 30 分）}

\begin{enumerate}[label=\textbf{\arabic*.}, itemsep=8pt]
  \item 不等式 $\dfrac{x - 1}{x + 2} > 0$ 的解集为 \underline{\hspace{2.5cm}}。
  \item 计算：$\log_2 8 + \log_3 \dfrac{1}{3} = $ \underline{\hspace{2.5cm}}。
  \item 抛物线 $y^2 = 8x$ 的焦点坐标为 \underline{\hspace{2.5cm}}。
  \item 二项式 $\left(x + \dfrac{2}{x}\right)^6$ 的展开式中常数项为 \underline{\hspace{2.5cm}}。
  \item 圆 $x^2 + y^2 - 4x + 2y + 1 = 0$ 的半径为 \underline{\hspace{2.5cm}}。
  \item 已知 $\tan\alpha = 2$，则 $\dfrac{\sin\alpha + \cos\alpha}{\sin\alpha - \cos\alpha} = $ \underline{\hspace{2.5cm}}。
\end{enumerate}

%======================================================================
\section*{三、解答题（本大题共 4 小题，共 80 分。解答须写出文字说明、证明过程或演算步骤）}

\begin{enumerate}[label=\textbf{\arabic*.}, itemsep=10pt]
  \item （本题满分 18 分）
  已知函数 $f(x) = 2\sqrt{3}\sin x\cos x + 2\cos^2 x - 1$。
  \begin{enumerate}[label=(\arabic*), itemsep=4pt]
    \item 求 $f(x)$ 的最小正周期；
    \item 求 $f(x)$ 在区间 $\left[0, \dfrac{\pi}{2}\right]$ 上的最大值和最小值。
  \end{enumerate}

  \item （本题满分 20 分）
  如图，在正方体 $ABCD\text{-}A_1B_1C_1D_1$ 中，$E$ 为棱 $CC_1$ 的中点。
  \begin{enumerate}[label=(\arabic*), itemsep=4pt]
    \item 求证：$BD \perp AC$；
    \item 求直线 $B_1E$ 与平面 $ABCD$ 所成角的正切值。
  \end{enumerate}

  \item （本题满分 20 分）
  设数列 $\{a_n\}$ 的前 $n$ 项和为 $S_n$，已知 $a_1 = 1$，且 $S_{n+1} = 2S_n + 1$。
  \begin{enumerate}[label=(\arabic*), itemsep=4pt]
    \item 证明数列 $\{a_n + 1\}$ 是等比数列；
    \item 求数列 $\{a_n\}$ 的通项公式与前 $n$ 项和。
  \end{enumerate}

  \item （本题满分 22 分）
  已知椭圆 $C\colon \dfrac{x^2}{a^2} + \dfrac{y^2}{b^2} = 1\ (a > b > 0)$ 的离心率为 $\dfrac{\sqrt{2}}{2}$，
  且经过点 $\left(\sqrt{2}, 1\right)$。
  \begin{enumerate}[label=(\arabic*), itemsep=4pt]
    \item 求椭圆 $C$ 的方程；
    \item 设直线 $l\colon y = x + m$ 与椭圆 $C$ 交于 $P$、$Q$ 两点，若 $\triangle OPQ$ 的面积为
    $\dfrac{\sqrt{2}}{2}$，求 $m$ 的值。
  \end{enumerate}
\end{enumerate}

\end{document}
